\section*{Chapitre 1 \textemdash\ Logique}
\addcontentsline{toc}{section}{Chapitre 1 \textemdash\ Logique}

%────────────────────────────────────────────────────────────────

\subsection*{Exercices}
\addcontentsline{toc}{subsection}{Exercices}

\solutiontitle{1}{1}{Propositions ou non}

\begin{enumerate}
  \item \textbf{Proposition vraie.}\enspace $5\times4=20$.
  \item \textbf{Proposition fausse.}\enspace $\sqrt{2}\approx1{,}414\notin\N$.
  \item \textbf{Non~: pr\'edicat.}\enspace La valeur de $x+3=7$ d\'epend de $x$.
  \item \textbf{Proposition vraie.}\enspace Tout triangle \'equilat\'eral a deux c\^ot\'es \'egaux.
  \item \textbf{Proposition vraie.}\enspace $0\times n=0$ pour tout $n$.
  \item \textbf{Non~: question.}\enspace Une question n'a pas de valeur de v\'erit\'e.
  \item \textbf{Proposition vraie.}\enspace $\pi\approx3{,}14159>3$.
  \item \textbf{Non~: pr\'edicat.}\enspace La valeur d\'epend de $n$.
  \item \textbf{Proposition vraie.}\enspace $\N\subset[0,+\infty)$ par d\'efinition.
  \item \textbf{Non~: ordre.}\enspace Un imp\'eratif n'est ni vrai ni faux.
\end{enumerate}

%────────────────────────────────────────────────────────────────
\solutiontitle{1}{2}{N\'egations simples}

\begin{enumerate}
  \item $17$ n'est \textbf{pas} un nombre premier. (Fausse~: $17$ est premier.)
  \item Le triangle $ABC$ n'est \textbf{pas} rectangle.
  \item $x < 5$.
  \item Il ne pleut pas \textbf{ou} il ne fait pas froid.
        (De Morgan~: $\lnot(P\land Q)\equiv\lnot P\lor\lnot Q$.)
  \item Il ne pleut pas \textbf{et} il ne fait pas chaud.
        ($\lnot(P\lor Q)\equiv\lnot P\land\lnot Q$.)
  \item $\exists n\in\N,\ n^2<n$. (Fausse~: $n(n-1)\geq0$ pour tout $n\in\N$.)
  \item $\forall x\in\R,\ x^2\ne-1$. (Vraie.)
  \item $\exists n\in\Z,\ 2n$ est impair. (Fausse.)
\end{enumerate}

%────────────────────────────────────────────────────────────────
\solutiontitle{1}{3}{Tables de v\'erit\'e}

\renewcommand{\arraystretch}{1.25}
\begin{center}
\begin{tabular}{|cc|cccccc|}
\hline
\rowcolor{BN!12}
$P$ & $Q$ & $\lnot P\lor Q$ & $P\land\lnot Q$ & $(P\lor Q)\land\lnot P$ &
$\lnot(P\land Q)$ & $\lnot P\land\lnot Q$ & $\lnot P\lor\lnot Q$ \\
\hline
\vrai & \vrai & \vrai & \faux & \faux & \faux & \faux & \faux \\
\vrai & \faux & \faux & \vrai & \faux & \vrai & \faux & \vrai \\
\faux & \vrai & \vrai & \faux & \vrai & \vrai & \faux & \vrai \\
\faux & \faux & \vrai & \faux & \faux & \vrai & \vrai & \vrai \\
\hline
\end{tabular}
\end{center}
Remarquer~: la formule~(4) $\lnot(P\land Q)$ est identique \`a~(6) $\lnot P\lor\lnot Q$~:
c'est la premi\`ere loi de De~Morgan.
De m\^eme,~(5) $\lnot P\land\lnot Q$ est la n\'egation de $P\lor Q$.

%────────────────────────────────────────────────────────────────
\solutiontitle{1}{4}{Valeur de v\'erit\'e d'une implication}

\begin{enumerate}
  \item \textbf{Vraie.}\enspace $P=\vrai$, $Q=\vrai$~: $P\Rightarrow Q$ est vraie.
  \item \textbf{Vraie.}\enspace $P=\faux$~: une implication \`a hypoth\`ese fausse est toujours vraie,
        quelle que soit la conclusion.
  \item \textbf{Fausse.}\enspace $P=\vrai$ et $Q=\faux$~: c'est l'unique cas o\`u une implication
        est fausse.
  \item \textbf{Vraie.}\enspace $P$~: \og$n$ divisible par~$4$\fg, $Q$~: \og$n$ divisible par~$2$\fg.
        Si $4\mid n$, alors $n=4k$, donc $n=2(2k)$~: $2\mid n$.
  \item \textbf{Fausse.}\enspace $n=2$ est divisible par $2$, mais pas par $4$.
\end{enumerate}

%────────────────────────────────────────────────────────────────
\solutiontitle{1}{5}{Quantificateurs~: traduction}

\begin{enumerate}
  \item \og Pour tout entier naturel $x$, $x+1>x$.\fg\enspace \textbf{Vraie.}
  \item \og Il existe un entier relatif dont le carr\'e vaut $-4$.\fg\enspace
        \textbf{Fausse}~: un carr\'e est toujours $\geq0$.
  \item \og Pour tout r\'eel $x$, $x^2\geq0$.\fg\enspace \textbf{Vraie.}
  \item \og Il existe un rationnel dont le carr\'e vaut $2$.\fg\enspace
        \textbf{Fausse.} ($\sqrt{2}\notin\Q$.)
  \item \og Pour tout r\'eel non nul, $1/x>0$.\fg\enspace
        \textbf{Fausse.} Contre-exemple~: $x=-1$.
  \item \og Il existe $n\in\N$ tel que $n^2=2n$.\fg\enspace
        \textbf{Vraie.} $n=0$ ou $n=2$ conviennent.
\end{enumerate}

%────────────────────────────────────────────────────────────────
\solutiontitle{1}{6}{\'Ecrire avec des quantificateurs}

\begin{enumerate}
  \item Pour tout quadrilat\`ere $Q$, si $Q$ est un carr\'e, alors $Q$ est un rectangle.
        En abr\'eg\'e~:
        \[\forall Q,\quad Q\text{ carr\'e}\Rightarrow Q\text{ rectangle}.\]
  \item $\exists n\in\N,\ (n \text{ est pair et } n \text{ est premier})$.
        (Vrai~: $n=2$.)
  \item $\forall a,b\in\Z,\ (2\mid a\ \text{et}\ 2\mid b)\Rightarrow 2\mid(a+b)$.
  \item $\forall x\in\R,\ \exists y\in\R,\ x+y=0$. (On prend $y=-x$.)
  \item $\lnot\bigl(\exists n\in\Z,\ 0<n<1\bigr)$, autrement dit
        $\forall n\in\Z,\ n\leq0\ \text{ou}\ n\geq1$.
\end{enumerate}

%────────────────────────────────────────────────────────────────
\solutiontitle{1}{7}{N\'egation avec quantificateurs}

\begin{enumerate}
  \item $\exists x\in\R,\ x^2\leq0$. \textbf{Vraie}~: $x=0$ donne $x^2=0$.
        (La proposition originale est donc fausse.)
  \item $\forall n\in\N,\ 2n+1\ne0$. \textbf{Vraie.}
        (L'originale est fausse~: $2n+1\geq1>0$.)
  \item $\exists n\in\N,\ n^2<n$. \textbf{Fausse.}
        (L'originale est vraie~: $n^2-n=n(n-1)\geq0$.)
  \item $\forall x\in\R,\ x^3\ne-8$. \textbf{Fausse.}
        (L'originale est vraie~: $x=-2$ convient.)
  \item $\exists x\in\R,\ \exists y\in\R,\ x+y\ne y+x$. \textbf{Fausse.}
        (L'originale est vraie~: l'addition est commutative dans $\R$.
        La n\'egation est donc fausse.)
\end{enumerate}

%────────────────────────────────────────────────────────────────
\solutiontitle{1}{8}{R\'eciproque et contrapos\'ee}

\begin{enumerate}
  \item Implication~: si $ABC$ est \'equilat\'eral, alors il est isoc\`ele.
        \textbf{Vraie.} Contrapos\'ee~: si $ABC$ n'est pas isoc\`ele, il n'est pas \'equilat\'eral.
        R\'eciproque~: si $ABC$ est isoc\`ele, il est \'equilat\'eral. \textbf{Fausse}
        (un triangle 5-5-4 est isoc\`ele non \'equilat\'eral).
  \item Implication~: $x=3\Rightarrow x^2=9$. \textbf{Vraie.}
        Contrapos\'ee~: $x^2\ne9\Rightarrow x\ne3$. Vraie.
        R\'eciproque~: $x^2=9\Rightarrow x=3$. \textbf{Fausse}~: $x=-3$ convient.
  \item Implication~: $6\mid n\Rightarrow3\mid n$. \textbf{Vraie.}
        Contrapos\'ee~: $3\nmid n\Rightarrow6\nmid n$. Vraie.
        R\'eciproque~: $3\mid n\Rightarrow6\mid n$. \textbf{Fausse}~: $n=3$ convient.
  \item Implication~: $x>1\Rightarrow x^2>1$. \textbf{Vraie.}
        Contrapos\'ee~: $x^2\leq1\Rightarrow x\leq1$. Vraie.
        R\'eciproque~: $x^2>1\Rightarrow x>1$. \textbf{Fausse}~: $x=-2$ convient.
\end{enumerate}

%────────────────────────────────────────────────────────────────
\solutiontitle{1}{9}{Contre-exemples}

\begin{enumerate}
  \item $n=0$~: $0^2=0\not>0$. (Ou $n=1$~: $1=1\not>1$.)
  \item $x=-3$~: $\sqrt{(-3)^2}=\sqrt{9}=3\ne-3$.
  \item $x=-2$~: $(-2)^2=4$ mais $x\ne2$.
  \item $n=4$~: $4$ est un multiple de $4$, mais $4/8=0{,}5\notin\Z$.
  \item $n=-1$~: $|-1|=1\ne-1$.
\end{enumerate}

%────────────────────────────────────────────────────────────────
\solutiontitle{1}{10}{Implication ou \'equivalence~?}

\begin{enumerate}
  \item $x=2\Rightarrow x^2=4$~: \textbf{vraie.}
        $x^2=4\Rightarrow x=2$~: \textbf{fausse} (contre-exemple $x=-2$).
        Conclusion~: $x=2\Rightarrow x^2=4$ mais non $\Leftrightarrow$.
  \item $n$ pair $\Rightarrow n^2$ pair~: vraie ($n=2k\Rightarrow n^2=4k^2$).
        $n^2$ pair $\Rightarrow n$ pair~: vraie (par contrapos\'ee~: $n$ impair $\Rightarrow n^2$ impair).
        Conclusion~: \textbf{\'equivalence} $n$ pair $\Leftrightarrow n^2$ pair.
  \item $6\mid n\Rightarrow 2\mid n$~: vraie.
        $2\mid n\Rightarrow 6\mid n$~: fausse (contre-exemple~: $n=2$, $6\nmid2$).
        Conclusion~: $6\mid n\Rightarrow 2\mid n$ seulement.
  \item $x=3\Rightarrow|x|=3$~: vraie.
        $|x|=3\Rightarrow x=3$~: fausse ($x=-3$ convient).
        Conclusion~: $x=3\Rightarrow|x|=3$ seulement.
  \item $x=0\Rightarrow xy=0$~: vraie.
        $xy=0\Rightarrow x=0$~: fausse (prendre $x=1$, $y=0$).
        Conclusion~: $x=0\Rightarrow xy=0$ seulement.
\end{enumerate}

%────────────────────────────────────────────────────────────────
\solutiontitle{1}{11}{Table~: contrapos\'ee}

\renewcommand{\arraystretch}{1.35}
\begin{center}
\begin{tabular}{|cc|cc|cc|}
\hline
\rowcolor{BN!12}
$P$ & $Q$ & $\lnot Q$ & $\lnot P$
    & $P\Rightarrow Q$ & $\lnot Q\Rightarrow\lnot P$ \\
\hline
\vrai & \vrai & \faux & \faux & \vrai & \vrai \\
\vrai & \faux & \vrai & \faux & \faux & \faux \\
\faux & \vrai & \faux & \vrai & \vrai & \vrai \\
\faux & \faux & \vrai & \vrai & \vrai & \vrai \\
\hline
\end{tabular}
\end{center}
Les colonnes 5 et 6 sont identiques~: $P\Rightarrow Q\equiv\lnot Q\Rightarrow\lnot P$.\hfill$\square$

%────────────────────────────────────────────────────────────────
\solutiontitle{1}{12}{Table~: connecteur $\Leftrightarrow$}

\begin{center}
\renewcommand{\arraystretch}{1.35}
\begin{tabular}{|cc|cc|c|c|}
\hline
\rowcolor{BN!12}
$P$ & $Q$ & $P\Rightarrow Q$ & $Q\Rightarrow P$
    & $(P\Rightarrow Q)\land(Q\Rightarrow P)$ & $P\Leftrightarrow Q$ \\
\hline
\vrai & \vrai & \vrai & \vrai & \vrai & \vrai \\
\vrai & \faux & \faux & \vrai & \faux & \faux \\
\faux & \vrai & \vrai & \faux & \faux & \faux \\
\faux & \faux & \vrai & \vrai & \vrai & \vrai \\
\hline
\end{tabular}
\end{center}
Les colonnes 5 et 6 co\"incident~: $P\Leftrightarrow Q\equiv(P\Rightarrow Q)\land(Q\Rightarrow P)$.\hfill$\square$

%────────────────────────────────────────────────────────────────
\solutiontitle{1}{13}{Calcul avec valeurs donn\'ees}

$P=\vrai$, $Q=\faux$, $R=\vrai$.
\begin{multicols}{2}
\begin{enumerate}
  \item $P\land Q = \vrai\land\faux = \faux$
  \item $P\lor Q = \vrai\lor\faux = \vrai$
  \item $P\Rightarrow Q = \vrai\Rightarrow\faux = \faux$
  \item $Q\Rightarrow P = \faux\Rightarrow\vrai = \vrai$
  \item $(P\lor Q)\Rightarrow R = \vrai\Rightarrow\vrai = \vrai$
  \item $\lnot P\lor(Q\land R)=\faux\lor(\faux\land\vrai)=\faux\lor\faux=\faux$
  \item $P\Leftrightarrow R = \vrai\Leftrightarrow\vrai = \vrai$
  \item $\lnot(P\land Q)\lor R = \vrai\lor\vrai = \vrai$
\end{enumerate}
\end{multicols}

%────────────────────────────────────────────────────────────────
\solutiontitle{1}{14}{Forme logique de r\'esultats math\'ematiques}

\begin{enumerate}
  \item $\forall n\in\Z,\ (n \text{ impair})\Rightarrow(n^2 \text{ impair})$.
  \item $\forall$ triangle $ABC$~:
        $(\text{rectangle en }C)\Leftrightarrow(AB^2=AC^2+BC^2)$.
  \item $\forall n\in\Z,\ \bigl((2\mid n)\land(3\mid n)\bigr)\Rightarrow6\mid n$.
  \item $\forall k\in\Z,\ 2\mid(4k+2)$.
\end{enumerate}

%────────────────────────────────────────────────────────────────
\solutiontitle{1}{15}{Priorit\'e des connecteurs}

Les priorit\'es~: $\lnot$ puis $\land$ puis $\lor$ puis $\Rightarrow$.

\begin{enumerate}
  \item $\lnot P\land Q\lor R=((\lnot P)\land Q)\lor R$.
        Elle est vraie ssi $R$ est vraie, ou si $P$ est fausse et $Q$ vraie.
  \item $P\land Q\Rightarrow R=(P\land Q)\Rightarrow R$.
        Elle est fausse seulement quand $P$ et $Q$ sont vraies et $R$ fausse.
  \item $P\lor Q\land\lnot R=P\lor\bigl(Q\land(\lnot R)\bigr)$.
        Elle est fausse seulement quand $P$ est fausse et que $Q\land\lnot R$ est fausse.
\end{enumerate}

%────────────────────────────────────────────────────────────────
\solutiontitle{1}{16}{De Morgan~: simplifications}

\begin{enumerate}
  \item $\lnot(\lnot P\land Q)\equiv\lnot(\lnot P)\lor\lnot Q\equiv P\lor\lnot Q$.
  \item $\lnot(P\lor\lnot Q)\equiv\lnot P\land Q$.
  \item $\lnot\bigl((P\Rightarrow Q)\land P\bigr)
        \equiv\lnot(P\Rightarrow Q)\lor\lnot P
        \equiv(P\land\lnot Q)\lor\lnot P$.
  \item $\lnot(\lnot P\lor\lnot Q)\equiv P\land Q$.
  \item $\lnot\bigl(\lnot(P\land Q)\bigr)\equiv P\land Q$.
  \item $\lnot(P\Leftrightarrow Q)
        \equiv(P\land\lnot Q)\lor(Q\land\lnot P)$ (ou exclusif).
\end{enumerate}

%────────────────────────────────────────────────────────────────
\solutiontitle{1}{17}{Table \`a trois variables}

\renewcommand{\arraystretch}{1.2}
\begin{center}
\small
\begin{tabular}{|ccc|c|c|c|c|}
\hline
\rowcolor{BN!12}
$P$&$Q$&$R$ & $P\Rightarrow Q$ & $Q\Rightarrow R$ &
$(P\Rightarrow Q)\land(Q\Rightarrow R)$ & $P\Rightarrow R$ \\
\hline
\vrai&\vrai&\vrai & \vrai & \vrai & \vrai & \vrai \\
\vrai&\vrai&\faux & \vrai & \faux & \faux & \faux \\
\vrai&\faux&\vrai & \faux & \vrai & \faux & \vrai \\
\vrai&\faux&\faux & \faux & \vrai & \faux & \faux \\
\faux&\vrai&\vrai & \vrai & \vrai & \vrai & \vrai \\
\faux&\vrai&\faux & \vrai & \faux & \faux & \vrai \\
\faux&\faux&\vrai & \vrai & \vrai & \vrai & \vrai \\
\faux&\faux&\faux & \vrai & \vrai & \vrai & \vrai \\
\hline
\end{tabular}
\end{center}
Quand $(P\Rightarrow Q)\land(Q\Rightarrow R)$ est vraie, $P\Rightarrow R$ l'est aussi.
La r\'eciproque est fausse (ligne 6). Cela confirme la \textbf{transitivit\'e de l'implication}~:
si $P\Rightarrow Q$ et $Q\Rightarrow R$, alors $P\Rightarrow R$.\hfill$\square$

%────────────────────────────────────────────────────────────────
\solutiontitle{1}{18}{Tautologies et contradictions}

\begin{enumerate}
  \item $P\lor\lnot P$~: $P=\vrai\Rightarrow\vrai\lor\faux=\vrai$~;
        $P=\faux\Rightarrow\faux\lor\vrai=\vrai$. \textbf{Tautologie.}
  \item $P\land\lnot P$~: toujours \faux. \textbf{Contradiction.}
  \item $P\Rightarrow(Q\Rightarrow P)$~: si $P=\faux$, c'est vrai.
        Si $P=\vrai$, alors $Q\Rightarrow P=Q\Rightarrow\vrai=\vrai$.
        \textbf{Tautologie.}
  \item $(P\Rightarrow Q)\Leftrightarrow(\lnot P\lor Q)$~:
        $P\Rightarrow Q$ est vraie ssi $P=\faux$ ou $Q=\vrai$,
        exactement quand $\lnot P\lor Q$ est vraie. \textbf{Tautologie.}
\end{enumerate}

%────────────────────────────────────────────────────────────────
\solutiontitle{1}{19}{Divisibilit\'e et implication}

$P$~: $4\mid n$~; $Q$~: $2\mid n$.

\begin{enumerate}
  \item $P\Rightarrow Q$~: si $n=4k$ alors $n=2(2k)$~: \textbf{vraie.}
        $Q\Rightarrow P$~: fausse. Contre-exemple~: $n=2$, $2\mid2$ mais $4\nmid2$.
  \item $P\Leftrightarrow Q$ est fausse (les sens ne sont pas \'equivalents).
  \item Contrapos\'ee de $P\Rightarrow Q$~: si $2\nmid n$ alors $4\nmid n$.
        V\'erification directe~: si $n$ est impair, alors $n\ne4k$ pour tout entier $k$.
  \item Si $a\mid b$ et $b\mid c$, alors $a\mid c$ (\textbf{transitivit\'e}).
        Preuve~: $b=ka$ et $c=\ell b$, donc $c=\ell ka=(k\ell)a$, d'o\`u $a\mid c$.\hfill$\square$
\end{enumerate}

%────────────────────────────────────────────────────────────────
\solutiontitle{1}{20}{Quantificateurs imbriqu\'es}

\begin{enumerate}
  \item \og Pour tout r\'eel $x$, il existe un r\'eel $y$ tel que $y=x+1$.\fg\enspace
        \textbf{Vraie.} Prendre $y=x+1$.
        N\'egation~: $\exists x\in\R,\ \forall y\in\R,\ y\ne x+1$.
  \item \og Il existe un r\'eel $x$ tel que pour tout r\'eel $y$, $y=x+1$.\fg\enspace
        \textbf{Fausse.} $x+1$ ne peut pas \'egaler tous les r\'eels simultanément.
        N\'egation~: $\forall x\in\R,\ \exists y\in\R,\ y\ne x+1$. (Vraie.)
  \item \og Pour tout r\'eel $x$, il existe $y$ tel que $x+y=0$.\fg\enspace
        \textbf{Vraie.} Prendre $y=-x$.
        N\'egation~: $\exists x,\ \forall y,\ x+y\ne0$.
  \item \og Il existe un r\'eel $x$ tel que pour tout r\'eel $y$, $x\cdot y=y$.\fg\enspace
        \textbf{Vraie.} $x=1$ convient~: $1\cdot y=y$ pour tout $y\in\R$.
        N\'egation~: $\forall x\in\R,\ \exists y\in\R,\ x\cdot y\ne y$. (Fausse.)
\end{enumerate}

%────────────────────────────────────────────────────────────────
\solutiontitle{1}{21}{Raisonnement par contrapos\'ee}

\textbf{(1)} Contrapos\'ee~: si $n$ est pair, alors $n^2$ est pair.
$n=2k\Rightarrow n^2=4k^2=2(2k^2)$.\hfill$\square$

\textbf{(2)} Contrapos\'ee~: si $n$ est impair, alors $n^2$ est impair.
$n=2k+1\Rightarrow n^2=4k^2+4k+1=2(2k^2+2k)+1$.\hfill$\square$

\textbf{(3)} $n^2+n=n(n+1)$. Si $n$ est impair, alors $n+1$ est pair, donc $n(n+1)$ est pair.
La contrapos\'ee de \og$n^2+n$ impair $\Rightarrow n$ pair\fg\ est \og$n$ impair $\Rightarrow n^2+n$ pair\fg, qui est vraie.\hfill$\square$

\textbf{(4)} Contrapos\'ee de $ab\ne0\Rightarrow(a\ne0$ et $b\ne0)$~:
si $a=0$ ou $b=0$, alors $ab=0$. Trivial.\hfill$\square$

%────────────────────────────────────────────────────────────────
\solutiontitle{1}{22}{G\'eom\'etrie et logique}

$P$~: \og$ABC$ isoc\`ele en $A$\fg~; $Q$~: \og$\hat B=\hat C$\fg.

\begin{enumerate}
  \item $P\Rightarrow Q$~: si $AB=AC$, alors $\hat B=\hat C$ (propri\'et\'e du triangle isoc\`ele).
        \textbf{Vraie.}
  \item $Q\Rightarrow P$~: si $\hat B=\hat C$, alors par la r\'eciproque du th\'eor\`eme
        de l'isosc\`ele (d\'emontr\'e par la somme des angles et la loi des sinus),
        $AB=AC$. \textbf{Vraie.}
        Donc $P\Leftrightarrow Q$.
  \item \textbf{Contrapos\'ee de $P\Rightarrow Q$~:}
        $\lnot Q\Rightarrow\lnot P$, soit~:
        \og Si $\hat B\ne\hat C$, alors le triangle n'est pas isoc\`ele en $A$\fg.
        D\'emonstration~: si $\hat B\ne\hat C$, supposons par l'absurde $AB=AC$.
        Alors les angles de la base seraient \'egaux (propri\'et\'e de l'isoc\`ele),
        contradiction. Donc $AB\ne AC$~: le triangle n'est pas isoc\`ele en $A$.\hfill$\square$
  \item $\lnot P$~: \og $AB\ne AC$\fg\ (triangle non isoc\`ele en $A$).
        $\lnot Q$~: \og $\hat B\ne\hat C$\fg.
\end{enumerate}

%────────────────────────────────────────────────────────────────
\solutiontitle{1}{23}{N\'egation de propositions complexes}

\begin{enumerate}
  \item N\'egation~: \og Il existe un entier pair qui n'est ni divisible par $4$
        ni divisible par $6$.\fg
        En symboles~: $\exists n\in\Z,\ (2\mid n)\ \land\ (4\nmid n)\ \land\ (6\nmid n)$.
        (Vraie~: $n=2$ est pair, non divisible par $4$ et non divisible par $6$.)
  \item N\'egation~: \og Pour tout entier $n$, $n^2\geq n$ ou $n\geq n^3$.\fg
        En symboles~: $\forall n\in\Z,\ n^2\geq n\ \lor\ n\geq n^3$.
  \item N\'egation~: \og Il existe un r\'eel $x>0$ tel que, pour tout $y>0$, $y\geq x$.\fg
        En symboles~: $\exists x>0,\ \forall y>0,\ y\geq x$.
  \item N\'egation~: \og Il existe un nombre premier qui n'est ni impair ni
        \`egal \`a $2$.\fg\ Cette proposition est fausse~: le seul nombre premier
        pair est $2$.
\end{enumerate}

%────────────────────────────────────────────────────────────────
\solutiontitle{1}{24}{Erreurs classiques de raisonnement}

\begin{enumerate}
  \item \textbf{Erreur~: affirmation du cons\'equent.}
        Les chats et les chiens sont des animaux, mais cela ne les rend pas identiques.
        La forme correcte serait~: si $P\Rightarrow R$ et $Q\Rightarrow R$, on ne peut pas
        conclure $P\Rightarrow Q$.
  \item \textbf{Erreur~: l'implication $x^2=4\Rightarrow x=2$ est fausse.}
        $x^2=4$ admet deux solutions~: $x=2$ et $x=-2$. On ne peut pas \'ecrire
        \og si $x^2=4$ alors $x=2$\fg.
  \item \textbf{Erreur~: g\'en\'eralisation h\^ative.}
        La v\'erification d'un nombre fini de cas ne prouve pas une assertion universelle.
        Avec $n=40$, on obtient $40^2+40+41=41^2$, qui n'est pas premier.
  \item \textbf{Erreur~: division par z\'ero.}
        Puisque $a=b$, on a $a-b=0$~: le passage de
        $(a-b)(a+b)=b(a-b)$ \`a $a+b=b$ est interdit.
\end{enumerate}

%────────────────────────────────────────────────────────────────
\solutiontitle{1}{25}{Double implication~: $n$ pair $\Leftrightarrow$ $n^2$ pair}

\textbf{Sens direct ($\Rightarrow$)~:} Si $n=2k$, alors $n^2=4k^2=2(2k^2)$~: pair.\hfill$\square$

\textbf{Sens r\'eciproque ($\Leftarrow$)~:} Par contrapos\'ee~: si $n$ est impair,
$n=2k+1$, donc $n^2=4k^2+4k+1=2(2k^2+2k)+1$~: impair.\hfill$\square$

Les deux sens \'etant d\'emontr\'es~: $n$ pair $\Leftrightarrow$ $n^2$ pair.

%────────────────────────────────────────────────────────────────
\solutiontitle{1}{26}{Conditions n\'ecessaires et suffisantes}

\begin{enumerate}
  \item $6\mid n$ pour $2\mid n$~: \textbf{condition suffisante.}
        $6\mid n\Rightarrow2\mid n$ (car $6=2\times3$).
        R\'eciproque fausse~: $n=2$, $2\mid2$ mais $6\nmid2$.
  \item $x=0$ pour $xy=0$~: \textbf{condition suffisante.}
        $x=0\Rightarrow xy=0$. R\'eciproque fausse~: $y=0$ suffit aussi.
  \item Un carr\'e pour un rectangle~: \textbf{condition suffisante},
        mais pas n\'ecessaire~: un rectangle de c\^ot\'es in\'egaux n'est pas un carr\'e.
  \item $\Delta=0$ pour une racine double~: \textbf{condition n\'ecessaire
        et suffisante} pour un trin\^ome du second degr\'e ($a\ne0$).
  \item $n$ impair pour $n^2$ impair~: \textbf{condition n\'ecessaire et suffisante}.
        Le sens direct est imm\'ediat~; si $n$ est pair, son carr\'e l'est aussi.
\end{enumerate}

%────────────────────────────────────────────────────────────────
\solutiontitle{1}{27}{D\'emonstration directe}

\begin{enumerate}
  \item $n=2k+1\Rightarrow n^3=(2k+1)^3=8k^3+12k^2+6k+1=2(4k^3+6k^2+3k)+1$~: impair.\hfill$\square$
  \item $n=2k+1\Rightarrow n^2+n=(2k+1)^2+(2k+1)=4k^2+4k+1+2k+1=4k^2+6k+2=2(2k^2+3k+1)$~: pair.\hfill$\square$
  \item $a=2p+1$, $b=2q+1\Rightarrow a+b=2p+2q+2=2(p+q+1)$~: pair.\hfill$\square$
  \item $a\mid b$ et $b\mid c$ signifient qu'il existe $k,\ell\in\Z$ tels que
        $b=ka$ et $c=\ell b$. Alors $c=\ell ka=(\ell k)a$,
        avec $\ell k\in\Z$. Ainsi $a\mid c$.\hfill$\square$
\end{enumerate}

%────────────────────────────────────────────────────────────────
\solutiontitle{1}{28}{Trouver l'erreur dans une preuve}

\begin{enumerate}
  \item \textbf{Erreur~: division par $n$ invalide pour $n=0$.}
        La division par $n$ n'est licite que si $n\ne0$.
        Or $n=0$ est un entier~: $0^2=0\geq0$ est vrai,
        mais l'\'etape \og diviser par $n$\fg\ est interdite pour $n=0$
        et ne prouve rien pour les autres valeurs.
        Correction~: $n^2\geq n\Leftrightarrow n^2-n\geq0\Leftrightarrow n(n-1)\geq0$.
        Vrai si $n\leq0$ ou $n\geq1$, donc vrai pour tout entier.
        Mais faux pour les r\'eels~: $n=1/2$ donne $1/4<1/2$.
        La proposition est vraie pour les entiers, mais la \emph{preuve} reste invalide.
  \item \textbf{Erreur~: l'implication \og carr\'e des membres\fg\ n'est pas une \'equivalence.}
        De $a^2+b^2=(a+b)^2$ on d\'eduit $2ab=0$, donc $a=0$ ou $b=0$.
        Cela prouve~: $\sqrt{a^2+b^2}=a+b$ implique $a=0$ ou $b=0$.
        Mais la conclusion de la \og preuve\fg\ est invers\'ee~:
        elle affirme que $\sqrt{a^2+b^2}=a+b$ \emph{en g\'en\'eral}.
        Contre-exemple~: $a=3, b=4$~: $\sqrt{9+16}=5\ne7$.
\end{enumerate}

%────────────────────────────────────────────────────────────────
\solutiontitle{1}{29}{Irrationnalit\'e de $\sqrt{2}$}

\textbf{Th\'eor\`eme.} $\sqrt{2}\notin\Q$.

\textbf{D\'emonstration} (par l'absurde). Supposons $\sqrt{2}=p/q$ avec $\pgcd(p,q)=1$.
Alors $p^2=2q^2$, donc $2\mid p^2$, et par l'exercice~21(2)~: $2\mid p$.
Posons $p=2m$. Alors $4m^2=2q^2$, d'o\`u $q^2=2m^2$, donc $2\mid q$.
Mais $2\mid\pgcd(p,q)$~: contradiction.\hfill$\square$

\textit{G\'en\'eralisation.} Si $k$ n'est pas un carr\'e parfait,
l'argument analogue (avec tout diviseur premier de $k$) m\`ene \`a la m\^eme contradiction.

%────────────────────────────────────────────────────────────────
\solutiontitle{1}{30}{Irrationnalit\'e de $\sqrt{3}$}

\textbf{Lemme.} $3\mid n^2\Rightarrow3\mid n$.
\emph{Preuve.} Par contrapos\'ee~: si $3\nmid n$, alors $n=3k+1$ ou $n=3k+2$.
Dans les deux cas, $n^2=9k^2+6k+1$ ou $n^2=9k^2+12k+4$, soit $n^2=3(\ldots)+1$ ou $3(\ldots)+4$~:
dans les deux cas $3\nmid n^2$.\hfill$\square$

\textbf{Th\'eor\`eme.} Supposons $\sqrt{3}=p/q$ avec $\pgcd(p,q)=1$.
Alors $p^2=3q^2$, donc $3\mid p^2$, et par le lemme $3\mid p$~: $p=3m$.
Alors $9m^2=3q^2$, d'o\`u $3\mid q$~: contradiction.\hfill$\square$

%────────────────────────────────────────────────────────────────
\solutiontitle{1}{31}{Suite et logique}

\begin{enumerate}
  \item $n=0$~: $0$~; $n=1$~: $0$~; $n=2$~: $2$~; $n=3$~: $6$~; $n=4$~: $12$~: tous pairs.\checkmark
  \item $n^2-n=n(n-1)$. Deux entiers cons\'ecutifs sont toujours l'un pair et l'autre impair,
        donc leur produit est pair.\hfill$\square$
  \item Si $n=2k$~: $n^2=4k^2=4k^2$, soit la forme $4k'$ avec $k'=k^2$.
        Si $n=2k+1$~: $n^2=4k^2+4k+1=4k(k+1)+1$, soit la forme $4k''+1$ avec $k''=k(k+1)$.
        Donc $n^2$ est de la forme $4k$ ou $4k+1$ pour tout $n\in\Z$.\hfill$\square$
\end{enumerate}

%────────────────────────────────────────────────────────────────
\solutiontitle{1}{32}{Infinit\'e des nombres premiers}

\textbf{Th\'eor\`eme} (Euclide). Il existe une infinit\'e de nombres premiers.

\textbf{D\'emonstration} (par l'absurde). Supposons qu'il n'en existe qu'un nombre fini
$p_1,\ldots,p_k$. Posons $N=p_1p_2\cdots p_k+1\geq2$.
$N$ admet un diviseur premier $p$. Mais pour tout $i$~:
$N = p_1p_2\cdots p_k+1$, donc $p_i\nmid N$ (reste $1$ dans la division)~: $p\notin\{p_1,\ldots,p_k\}$.
Contradiction.\hfill$\square$

%────────────────────────────────────────────────────────────────
\solutiontitle{1}{33}{Combinaisons logiques}

\begin{enumerate}
  \item \textbf{Exportation~:} $P\Rightarrow(Q\Rightarrow R)$~:
        si $P$ est vrai, il reste \`a montrer $Q\Rightarrow R$.
        $(P\land Q)\Rightarrow R$~: si $P$ et $Q$ sont vrais, montrer $R$.
        Les deux \'enonc\'es exigent la m\^eme chose~: v\'erifiable par table de v\'erit\'e (8 lignes).\hfill$\square$
  \item $(P\Rightarrow Q)\land(P\Rightarrow R)\equiv P\Rightarrow(Q\land R)$~:
        les deux affirment que si $P$ est vrai alors $Q$ et $R$ le sont tous les deux.
        Table de v\'erit\'e~: les deux formules co\"incident sur toutes les lignes.\hfill$\square$
  \item $P\lor(Q\land R)\equiv(P\lor Q)\land(P\lor R)$~:
        distribution de $\lor$ sur $\land$.
        V\'erifier par table~: si $P=\vrai$, les deux c\^ot\'es sont vrais.
        Si $P=\faux$~: gauche $=Q\land R$, droite $=Q\land R$.\hfill$\square$
\end{enumerate}

%────────────────────────────────────────────────────────────────
\solutiontitle{1}{34}{Raisonner sur les r\'eels}

\begin{enumerate}
  \item $x^2=(x-0)^2\geq0$ car un carr\'e est toujours $\geq0$.
        Donc $x^2+1\geq0+1=1>0$.\hfill$\square$
  \item Par l'absurde~: supposons qu'il existe un plus petit r\'eel $\varepsilon>0$.
        Alors $\varepsilon/2>0$ et $\varepsilon/2<\varepsilon$~: contradiction.\hfill$\square$
  \item $0\leq(x-y)^2=x^2-2xy+y^2$, donc $2xy\leq x^2+y^2$~:
        $xy\leq\dfrac{x^2+y^2}{2}$.\hfill$\square$
\end{enumerate}

%────────────────────────────────────────────────────────────────
\solutiontitle{1}{35}{Nombres rationnels et puissances}

\begin{enumerate}
  \item $2^n$ est pair (car $2\mid2^n$) et $3^m$ est impair (car $3$ est impair,
        et une puissance d'impair est impaire). Un nombre ne peut \^etre
        \`a la fois pair et impair~: $2^n\ne3^m$.\hfill$\square$
  \item Si $(2/3)^n=2^n/3^n$ \'etait entier, cela signifierait $3^n\mid2^n$~:
        impossible car $\pgcd(2,3)=1$ (ils n'ont aucun facteur premier commun).\hfill$\square$
  \item $(2/3)^n=2^n/3^n$. On a $2^n<3^n$ pour $n\geq1$ (car $2<3$, et la
        comparaison se conserve par puissance). Donc $0<(2/3)^n<1$, et
        $(2/3)^n$ est de plus en plus petit~: pour tout $\varepsilon>0$,
        choisir $n>\log\varepsilon/\log(2/3)$ convient.
\end{enumerate}

%────────────────────────────────────────────────────────────────
\solutiontitle{1}{36}{R\'eciproque du th\'eor\`eme de Pythagore}

\begin{enumerate}
  \item \textbf{Contrapos\'ee~:} si $AB^2\ne AC^2+BC^2$, alors le triangle n'est pas
        rectangle en $C$.
  \item \textbf{R\'eciproque~:} si $AB^2=AC^2+BC^2$, alors le triangle est rectangle en $C$.
        Cette r\'eciproque est \textbf{vraie} (c'est le th\'eor\`eme de Pythagore r\'eciproque).
  \item $5^2+12^2=25+144=169=13^2$~: le triangle de c\^ot\'es $5$, $12$, $13$
        est rectangle. De m\^eme, $6^2+8^2=36+64=100=10^2$~: celui de c\^ot\'es
        $6$, $8$, $10$ est rectangle.
  \item \textbf{Vraie.}\enspace Contrapos\'ee de la r\'eciproque~: si le triangle
        n'est pas rectangle en $C$, alors $AB^2\ne AC^2+BC^2$.
\end{enumerate}

%────────────────────────────────────────────────────────────────
\subsection*{Probl\`emes}
\addcontentsline{toc}{subsection}{Probl\`emes}
%────────────────────────────────────────────────────────────────

\psolutiontitle{1}{1}{Circuit logique}

\renewcommand{\arraystretch}{1.25}
\begin{center}
\small
\begin{tabular}{|ccc|c|c|c|c|c|}
\hline
\rowcolor{BN!12}
$A$ & $B$ & $C$ & $B\lor C$ & $\lnot A$
  & $A\land(B\lor C)$ & $\lnot A\land B\land C$ & $\Phi$ \\
\hline
\vrai&\vrai&\vrai & \vrai & \faux & \vrai & \faux & \vrai \\
\vrai&\vrai&\faux & \vrai & \faux & \vrai & \faux & \vrai \\
\vrai&\faux&\vrai & \vrai & \faux & \vrai & \faux & \vrai \\
\vrai&\faux&\faux & \faux & \faux & \faux & \faux & \faux \\
\faux&\vrai&\vrai & \vrai & \vrai & \faux & \vrai & \vrai \\
\faux&\vrai&\faux & \vrai & \vrai & \faux & \faux & \faux \\
\faux&\faux&\vrai & \vrai & \vrai & \faux & \faux & \faux \\
\faux&\faux&\faux & \faux & \vrai & \faux & \faux & \faux \\
\hline
\end{tabular}
\end{center}

Le courant passe pour $(\vrai,\vrai,\vrai)$, $(\vrai,\vrai,\faux)$,
$(\vrai,\faux,\vrai)$, $(\faux,\vrai,\vrai)$~: exactement les cas o\`u au moins
deux interrupteurs sont ferm\'es. $\Phi\equiv(A\land B)\lor(A\land C)\lor(B\land C)$~:
c'est la \textbf{fonction majorit\'e}.\hfill$\square$

\psolutiontitle{1}{2}{Chevaliers et coquins}

$P$~: \og$X$ est chevalier\fg~; $Q$~: \og$Y$ est chevalier\fg.

$X$ d\'eclare \og$Y$ est un coquin\fg, soit $\lnot Q$.
\begin{itemize}
  \item Si $P$~: la d\'eclaration de $X$ est vraie, donc $\lnot Q$~: $Y$ est coquin.
  \item Si $\lnot P$~: $X$ ment, donc $Q$~: $Y$ est chevalier.
\end{itemize}
$Y$ d\'eclare \og$X$ et $Z$ sont du m\^eme type\fg~: $(P\Leftrightarrow Z)$ o\`u $Z$~: \og$Z$ est chevalier\fg.
\begin{itemize}
  \item \textbf{Cas $P$~:} $Q=\faux$, donc $Y$ ment~: la d\'eclaration de $Y$ est fausse,
        i.e., $X$ et $Z$ sont de types diff\'erents~: $\lnot(P\Leftrightarrow Z)$,
        donc $Z=\faux$ (puisque $P=\vrai$).
  \item \textbf{Cas $\lnot P$~:} $Q=\vrai$, donc $Y$ dit vrai~:
        $X$ et $Z$ sont du m\^eme type, et $P=\faux$, donc $Z=\faux$.
\end{itemize}
\textbf{Conclusion.}
Dans les deux cas $Z=\faux$ (coquin). $X$ peut \^etre chevalier ou coquin~:
les deux sc\'enarios sont logiquement coh\'erents (l'\'enigme est ind\'etermin\'ee
pour $X$, mais $Z$ est toujours coquin).

\psolutiontitle{1}{3}{Portes logiques universelles}

Une porte NAND calcule $\lnot(P\land Q)$.

\begin{enumerate}
  \item $\lnot P = \lnot(P\land P) = \text{NAND}(P,P)$.\hfill$\square$
  \item \begin{align*}
        P\land Q
        &= \lnot\bigl(\lnot(P\land Q)\bigr)\\
        &= \text{NAND}\bigl(\text{NAND}(P,Q),\text{NAND}(P,Q)\bigr).
        \end{align*}
  \item \begin{align*}
        P\lor Q
        &= \lnot(\lnot P\land\lnot Q)\\
        &= \text{NAND}\bigl(\text{NAND}(P,P),\text{NAND}(Q,Q)\bigr).
        \end{align*}
  \item Puisqu'on peut exprimer $\lnot$, $\land$ et $\lor$ uniquement avec NAND,
        et que tout connecteur logique s'exprime avec ces trois-l\`a,
        les portes NAND sont \textbf{universelles}.\hfill$\square$
        (Idem pour NOR~: $\lnot P=\text{NOR}(P,P)$, $P\lor Q=\lnot(\text{NOR}(P,Q))$,
        etc.)
\end{enumerate}

\psolutiontitle{1}{4}{\'Enigme logique~: qui a fait quoi~?}

Exactement une d\'eclaration est vraie. Testons les trois cas~:

\textbf{Cas $A$ coupable~:} $A$ dit \og ce n'est pas moi\fg~: \textbf{faux}.
$B$ dit \og c'est $A$\fg~: \textbf{vrai}.
$C$ dit \og ce n'est pas $A$\fg~: \textbf{faux}.
Une seule d\'eclaration vraie~: \textbf{coh\'erent.}

\textbf{Cas $B$ coupable~:} $A$ dit \og ce n'est pas moi\fg~: vrai.
$B$ dit \og c'est $A$\fg~: faux.
$C$ dit \og ce n'est pas $A$\fg~: vrai.
Deux d\'eclarations vraies~: \textbf{impossible.}

\textbf{Cas $C$ coupable~:} $A$~: vrai~; $B$~: faux~; $C$~: vrai.
Deux d\'eclarations vraies~: \textbf{impossible.}

\textbf{Conclusion.} \textbf{$A$ a commis le vol.}

\psolutiontitle{1}{5}{Le berger, le loup, la ch\`evre et le chou}

\'Etat~: $(B,L,C,X)$ o\`u $\vrai$ = c\^ot\'e d'arriv\'ee.
\'Etat initial~: $(\faux,\faux,\faux,\faux)$.
\'Etat final~: $(\vrai,\vrai,\vrai,\vrai)$.
Contraintes~: ne jamais avoir $L=C\ne B$ (loup mange ch\`evre)
ni $C=X\ne B$ (ch\`evre mange chou).

\textbf{Solution~:}
\begin{enumerate}
  \item Traverser avec la ch\`evre~: $(\vrai,\faux,\vrai,\faux)$.
  \item Revenir seul~: $(\faux,\faux,\vrai,\faux)$.
  \item Traverser avec le loup~: $(\vrai,\vrai,\vrai,\faux)$.
  \item Revenir avec la ch\`evre~: $(\faux,\vrai,\faux,\faux)$.
  \item Traverser avec le chou~: $(\vrai,\vrai,\faux,\vrai)$.
  \item Revenir seul~: $(\faux,\vrai,\faux,\vrai)$.
  \item Traverser avec la ch\`evre~: $(\vrai,\vrai,\vrai,\vrai)$.
\end{enumerate}
Solution alternative~: remplacer \og loup\fg\ par \og chou\fg\ aux \'etapes~3 et~4.
Il y a exactement deux solutions.

\psolutiontitle{1}{6}{Parit\'e des puissances}

\begin{enumerate}
  \item Si $n$ est pair~: $n=2k$, $n^3=8k^3=2(4k^3)$~: pair.
        Si $n$ est impair~: $n=2k+1$, $n^3=(2k+1)^3=2(\ldots)+1$~: impair.
        Donc $n$ et $n^3$ ont la m\^eme parit\'e.\hfill$\square$
  \item $n=2m+1$~: $n^2=(2m+1)^2=4m^2+4m+1=4m(m+1)+1$.
        Parmi $m$ et $m+1$, l'un est pair~: $m(m+1)=2k$.
        Donc $n^2=8k+1$.\hfill$\square$
  \item Si $4\mid(a^2+b^2)$, alors $a^2+b^2$ est de la forme $4\ell$.
        Or $a^2$ et $b^2$ sont chacun de la forme $4k$ ou $4k+1$
        (voir exercice~31). Leur somme vaut $4\ell$ seulement si
        les deux sont de la forme $4k$, c'est-\`a-dire $a$ et $b$ sont pairs.\hfill$\square$
\end{enumerate}

\psolutiontitle{1}{7}{Divisibilit\'e et parit\'e des carr\'es}

\begin{enumerate}
  \item $n=2k\Rightarrow n^2=4k^2$~: de la forme $4k'$.
        $n=2k+1\Rightarrow n^2=4k^2+4k+1=4k(k+1)+1$~: de la forme $4k'+1$.
        Donc $n^2$ est de la forme $4k$ ou $4k+1$.\hfill$\square$
  \item $a^2+b^2$ vaut $4k$, $4k+1$ ou $4k+2$ (jamais $4k+3$).
        Donc $a^2+b^2$ n'est jamais de la forme $4k+3$.
        Les entiers de cette forme sont~: $3,7,11,15,\ldots$~:
        aucun d'eux ne s'\'ecrit comme somme de deux carr\'es.\hfill$\square$
  \item $7=4\cdot1+3$~: de la forme $4k+3$, donc impossible comme somme de deux carr\'es.
        $10=9+1=3^2+1^2$~: possible. $13=4+9=2^2+3^2$~: possible.
        $15=4\cdot3+3$~: de la forme $4k+3$, impossible.
  \item Sommes de deux carr\'es dans $\{1,\ldots,20\}$~:
        $1,2,4,5,8,9,10,13,16,17,18,20$. Les impossibles~: $3,6,7,11,12,14,15,19$.
\end{enumerate}

\psolutiontitle{1}{8}{Logique dans les preuves g\'eom\'etriques}

\begin{enumerate}
  \item $P\Rightarrow Q$ ($\vrai$)~; $P\Rightarrow R$ ($\vrai$)~;
        $Q\Rightarrow S$ ($\vrai$)~; $R\Rightarrow S$ ($\vrai$)~; $P\Rightarrow S$ ($\vrai$).
        R\'eciproquement~: $Q\not\Rightarrow P$ (un rectangle non carr\'e)~;
        $R\not\Rightarrow P$ (un losange non carr\'e).
        Donc $P\Rightarrow Q$ et $P\Rightarrow R$ mais $Q\not\Leftrightarrow P$.
  \item $\lnot P\lor(\lnot Q\lor\lnot R)$~: \og si $ABCD$ n'est pas un carr\'e,
        il n'est ni un rectangle, ni un losange\fg~: \textbf{fausse}
        (un rectangle qui n'est pas un carr\'e est un contre-exemple).
  \item $P\Leftrightarrow (Q\land R)$~:
        un carr\'e est un rectangle ET un losange, et r\'eciproquement
        un rectangle qui est aussi un losange est un carr\'e. \textbf{Vraie.}\hfill$\square$
\end{enumerate}

\psolutiontitle{1}{9}{Sudoku $4\times4$ par d\'eduction logique}

Grille de d\'epart~:
\[
\begin{array}{|cc|cc|}
\hline
  & 2 & 3 &   \\
3 &   &   & 2 \\
\hline
  & 3 & 4 &   \\
4 &   &   & 3 \\
\hline
\end{array}
\]

\textbf{\'Etape~1.} Colonne~1 contient $\{3,4\}$~: manquent $1$ et $2$.
Ligne~1 contient $\{2,3\}$~: manquent $1$ et $4$.
$(1,1)$ doit \^etre dans $\{1,2\}\cap\{1,4\}=\{1\}$~: \textbf{$(1,1)=1$}.
Ligne~1 compl\`ete~: $(1,4)=4$.

\textbf{\'Etape~2.} Ligne~4 contient $\{4,3\}$~: manquent $1$ et $2$.
Colonne~4 contient $\{4,2,3\}$~: manque $1$.
$(4,4)\in\{1,2\}\cap\{1\}=\{1\}$~: \textbf{$(4,4)=1$}. Donc $(4,3)=2$.

\textbf{\'Etape~3.} Colonne~2 contient $\{2,3\}$~: manquent $1$ et $4$.
Ligne~2 contient $\{3,2\}$~: manquent $1$ et $4$.
Bloc haut-gauche contient $\{1,2,3\}$~: manque $4$.
$(2,2)\in\{1,4\}\cap\{1,4\}\cap\{4\}=\{4\}$~: \textbf{$(2,2)=4$}. Donc $(2,3)=1$.

\textbf{\'Etape~4.} Ligne~3 contient $\{3,4\}$~: manquent $1$ et $2$.
Colonne~1 contient $\{1,3,4\}$~: manque $2$.
$(3,1)=2$, $(3,4)=1$.

\textbf{Solution unique~:}
\[
\begin{array}{|cc|cc|}
\hline
\mathbf{1} & 2 & 3 & \mathbf{4} \\
3 & \mathbf{4} & \mathbf{1} & 2 \\
\hline
\mathbf{2} & 3 & 4 & \mathbf{1} \\
4 & \mathbf{1} & \mathbf{2} & 3 \\
\hline
\end{array}
\]

\begin{enumerate}
  \item Chaque \'etape utilise une implication du type~:
        \og La case $(i,j)$ ne peut pas contenir $k$ car $k$ est
        d\'ej\`a pr\'esent dans la ligne/colonne/bloc~: donc elle contient~$\ell$.\fg
  \item $4$ implications principales ont \'et\'e utilis\'ees (une par case inconnue).
  \item La solution est unique~: chaque \'etape \'etait forc\'ee.
\end{enumerate}

\psolutiontitle{1}{10}{Sommes et formules alg\'ebriques}

\begin{enumerate}
  \item V\'erification~: $n=1$~: $1=1$\checkmark~; $n=2$~: $3=3$\checkmark~;
        $n=3$~: $6=6$\checkmark~; $n=4$~: $10=10$\checkmark~; $n=5$~: $15=15$\checkmark.
  \item Deux m\'ethodes pour $n=5$~: $1+2+3+4+5=15$ et $5\cdot6/2=15$. \'Egaux.\checkmark
  \item \textbf{Algorithme~:}
\begin{algorithm}[H]
\caption{V\'erification de la formule de Gauss}
\KwIn{un entier $n\geq1$}
$S \leftarrow 0$\;
\Pour{$k\leftarrow1$ \KwTo $n$}{$S\leftarrow S+k$\;}
\Si{$S = n\times(n+1)/2$}{\Afficher \og Formule v\'erifi\'ee\fg\;}
\Sinon{\Afficher \og Erreur\fg\;}
\end{algorithm}
\end{enumerate}

\psolutiontitle{1}{11}{Codage et logique~: parit\'e}

\begin{enumerate}
  \item \texttt{1011}~: trois $1$~: nombre impair. Bit de parit\'e $= 1$
        (pour que le total soit pair~: $3+1=4$, pair).
        Message envoy\'e~: \texttt{10111}.
  \item \texttt{10110}~: quatre $1$~: pair. Pas d'erreur d\'etect\'ee.
  \item \texttt{10100}~: trois $1$~: impair. \textbf{Erreur d\'etect\'ee.}
  \item Si deux erreurs se produisent, la parit\'e reste inchang\'ee~:
        le syst\`eme \textbf{ne peut pas} d\'etecter deux erreurs simultan\'ees.
        La parit\'e ne d\'etecte que les erreurs en nombre \emph{impair}.
\end{enumerate}

\psolutiontitle{1}{12}{D\'eduction logique~: tournoi de foot}

6 matchs au total (6 paires parmi 4 \'equipes). $A$ gagne 3 matchs~: 9 points.
Points restants distribu\'es~: $6\times3-9=9$ points pour $\{B,C,D\}$
mais avec les nuls possibles.

$B$ a 6 points. Comme $B$ a perdu contre $A$, elle a gagn\'e ses deux autres
matchs. Ainsi $B$ bat $C$ et $D$. Les points de $C$ et $D$ proviennent alors
uniquement de leur confrontation directe. Pour qu'ils soient \`a \`egalit\'e,
ce match est nul.

\textbf{Classement final~:} $A=9$, $B=6$, $C=D=1$.\hfill$\square$

\psolutiontitle{1}{13}{Logique et calcul alg\'ebrique}

\begin{enumerate}
  \item N\'egation~: \og Il existe $x\in\R$ tel que $x^2<0$.\fg
        Fausse~: la proposition originale est vraie.
  \item Contrapos\'ee~: \og Si $x\ne0$, alors $x^2\ne0$.\fg
        Preuve~: si $x\ne0$, alors $x^2=x\cdot x\ne0$ (produit de deux non-nuls).
  \item $ax^2+bx+c=0$ a deux racines distinctes $\Leftrightarrow\Delta>0$.
        R\'eciproque~: $\Delta>0\Rightarrow$ deux racines distinctes. \textbf{Vraie.}
  \item \og$\exists x\in\R,\ x^2+1=0$\fg~: n\'egation~:
        \og$\forall x\in\R,\ x^2+1\ne0$.\fg La n\'egation est vraie
        car $x^2\geq0$, donc $x^2+1\geq1>0$.
\end{enumerate}

\psolutiontitle{1}{14}{Algorithme et divisibilit\'e}

\begin{enumerate}
  \item Oui~: $n\bmod3=0\Leftrightarrow3\mid n$ et $n\bmod9\ne0\Leftrightarrow9\nmid n$.
        Les deux notations sont \'equivalentes.
  \item N\'egation~: \og$3\nmid n$ ou $9\mid n$\fg.
  \item Pour $n=3$~: $3\bmod9=3\ne0$~: affich\'e.
        $n=6$~: affich\'e. $n=9$~: $9\bmod9=0$~: non affich\'e.
        $n=12$~: affich\'e. $n=18$~: non affich\'e. $n=21$~: affich\'e.
  \item Entre $1$ et $100$~: multiples de $3$~: $\lfloor100/3\rfloor=33$.
        Multiples de $9$~: $\lfloor100/9\rfloor=11$.
        Multiples de $3$ mais pas $9$~: $33-11=\mathbf{22}$ entiers.
\end{enumerate}

%────────────────────────────────────────────────────────────────
\subsection*{D\'efis}
\addcontentsline{toc}{subsection}{D\'efis}
%────────────────────────────────────────────────────────────────

\noindent\phantomsection
{\color{BO}\bfseries\small D\'efi~1~\textendash\ \textit{Le paradoxe du menteur}}
\par\smallskip

\begin{enumerate}
  \item Si elle est vraie~: elle affirme qu'elle est fausse, donc elle est fausse.
        Si elle est fausse~: elle affirme qu'elle est fausse, ce qui est vrai.
        Dans les deux cas on obtient une contradiction.
  \item Cette phrase n'est pas une proposition car elle ne peut avoir de valeur de
        v\'erit\'e d\'etermin\'ee~: elle est \textbf{auto-r\'ef\'erentielle}.
        Le cadre logique du cours exige que les propositions aient une valeur de
        v\'erit\'e bien d\'efinie.
  \item Autres exemples~: \og Cette phrase contient cinq mots.\fg\
        (Non paradoxale~: vraie.) \og Cette phrase est vraie.\fg\
        (Non paradoxale~: tautologie.)
        \og Cette phrase est fausse\fg\ est la version paradoxale.
\end{enumerate}

\noindent\phantomsection
{\color{BO}\bfseries\small D\'efi~2~\textendash\ \textit{Construire sa propre table de v\'erit\'e}}
\par\smallskip

\begin{enumerate}
  \item Table du \textbf{ou exclusif} $P\oplus Q$~:
        \begin{center}\renewcommand{\arraystretch}{1.2}
        \begin{tabular}{|cc|c|}
        \hline\rowcolor{BN!10}$P$&$Q$&$P\oplus Q$\\\hline
        \vrai&\vrai&\faux\\
        \vrai&\faux&\vrai\\
        \faux&\vrai&\vrai\\
        \faux&\faux&\faux\\\hline
        \end{tabular}
        \end{center}
  \item $P\oplus Q\equiv(P\lor Q)\land\lnot(P\land Q)$~:
        v\'erifier ligne par ligne~: vrai seulement quand exactement un des deux est vrai.\checkmark
  \item $P\oplus Q\equiv(P\Rightarrow\lnot Q)\land(\lnot P\Rightarrow Q)$.
        V\'erification~:
        $(P\Rightarrow\lnot Q)$ est fausse seulement si $P=Q=\vrai$~;
        $(\lnot P\Rightarrow Q)$ est fausse seulement si $P=Q=\faux$.
        La conjonction est vraie exactement pour les deux lignes centrales.\checkmark
  \item XOR est aussi \textbf{universel}~:
        $\lnot P=P\oplus\vrai$~; $P\land Q=(P\oplus Q\oplus(P\lor Q))$, etc.
\end{enumerate}

\noindent\phantomsection
{\color{BO}\bfseries\small D\'efi~3~\textendash\ \textit{\'Enigmes logiques~: chevaliers et coquins}}
\par\smallskip

\begin{enumerate}
  \item $A$ dit \og nous sommes tous les deux des coquins\fg.
        Si $A$ est chevalier~: la phrase est vraie, donc $A$ est coquin~: contradiction.
        Donc \textbf{$A$ est coquin} (il ment). La phrase est fausse, donc
        ils ne sont \emph{pas} tous les deux coquins~: \textbf{$B$ est chevalier}.
  \item $C$ dit \og Au moins l'un de nous est un coquin.\fg
        Si $C$ est chevalier~: vrai, donc il y a au moins un coquin. Possible.
        Si $C$ est coquin~: la phrase est fausse, donc tous les deux sont chevaliers~:
        contradiction ($C$ serait chevalier).
        Donc \textbf{$C$ est chevalier}, et il y a bien un coquin~: \textbf{$D$ est coquin}.
  \item $E$ dit \og je suis un coquin\fg.
        Si $E$ est chevalier~: vrai, donc $E$ est coquin~: contradiction.
        Si $E$ est coquin~: il ment, donc $E$ n'est pas coquin~: contradiction.
        \textbf{Impossible}~: aucun logicien ne peut faire cette d\'eclaration coh\'eremment.
\end{enumerate}

\noindent\phantomsection
{\color{BO}\bfseries\small D\'efi~4~\textendash\ \textit{Bijections et taille d'ensembles}}
\par\smallskip

\begin{enumerate}
  \item $\{1,2,3\}\to\{a,b,c\}$~: bijection $1\mapsto a$, $2\mapsto b$, $3\mapsto c$. \textbf{Oui.}
  \item Bijection $\N\to\Z$~: $0\mapsto0$, $1\mapsto1$, $2\mapsto-1$, $3\mapsto2$, $4\mapsto-2$, \ldots
        En g\'en\'eral~: $2k\mapsto k$ et $2k+1\mapsto-k$.
        Donc $|\N|=|\Z|$~: \textbf{\'equipotents.}
  \item Bijection $\N\to\N\times\N$~: parcourir les diagonales~:
        $(0,0),(1,0),(0,1),(2,0),(1,1),(0,2),\ldots$ (argument de Cantor).
        Donc $|\N|=|\N\times\N|$.
  \item L'ensemble $\{0,1\}^\N$ (suites infinies de~0 et~1) est non d\'enombrable
        (preuve diagonale de Cantor)~: il est \textbf{strictement plus grand} que $\N$.
\end{enumerate}

\noindent\phantomsection
{\color{BO}\bfseries\small D\'efi~5~\textendash\ \textit{Raisonnement par l'absurde~: deux preuves classiques}}
\par\smallskip

\begin{enumerate}
  \item Supposition~: $7k$ est le plus grand multiple de $7$.
        Alors $7(k+1)=7k+7$ est aussi un multiple de $7$ et $7(k+1)>7k$~: contradiction.
        Donc il y a une infinit\'e de multiples de $7$.\hfill$\square$
  \item Supposition~: il existe un plus grand entier $N$.
        Alors $N+1$ est un entier et $N+1>N$~: contradiction.\hfill$\square$
  \item Supposition~: $\log_2 3=p/q$ avec $p,q\in\N^*$.
        Alors $2^p=3^q$. Mais $2^p$ est pair et $3^q$ est impair~: impossible.\hfill$\square$
\end{enumerate}

\noindent\phantomsection
{\color{BO}\bfseries\small D\'efi~6~\textendash\ \textit{Construire un syst\`eme de r\`egles logiques}}
\par\smallskip

\begin{enumerate}
  \item Applications du modus ponens~:
        \begin{itemize}
          \item \og Il pleut\fg\ + \og S'il pleut, le sol est mouill\'e\fg\ $\Rightarrow$
                \og Le sol est mouill\'e.\fg
          \item \og $n$ est pair\fg\ + \og Si $n$ est pair, $n^2$ est pair\fg\ $\Rightarrow$
                \og $n^2$ est pair.\fg
        \end{itemize}
  \item Le modus tollens dit~: de $P\Rightarrow Q$ et $\lnot Q$, conclure $\lnot P$.
        Cela correspond \`a la contrapos\'ee~: c'est la m\^eme r\`egle vue diff\'eremment.
        \textbf{Exemple~:} \og S'il a plu, le sol est mouill\'e.\fg\ \og Le sol est sec.\fg\
        $\Rightarrow$ \og Il n'a pas plu.\fg
  \item \textbf{D\'eduction en cha\^ine~:}
        $P\Rightarrow Q$, $Q\Rightarrow R$, $P$ est vrai $\Rightarrow$ $R$ est vrai.
        Par transitivit\'e (exercice~17), $P\Rightarrow R$, puis modus ponens.\hfill$\square$
  \item Le \textbf{raisonnement par cas} est la r\`egle~:
        de $(P\Rightarrow R)\land(Q\Rightarrow R)\land(P\lor Q)$, conclure $R$.
        Il est valide~: si $P\lor Q$ est vraie, l'un des deux cas s'applique,
        et dans les deux cas $R$ est vrai.
\end{enumerate}



\needspace{3\baselineskip}
